% fixture: amsmath/align-tag
% tier: T3
% engine: document
% covers: \tag and \tag* in align and align* (\tag@in@align: \df@tag is \tagform@{…} or \maketag@@@{…} in \normalfont, no counter stepped), a tag with a formula and one ending in an italic letter (\@@italiccorr), a \label on a \tag row reading the tag's text, \notag beside them, and the numbering going on after
\documentclass{article}
\usepackage{amsmath}
\begin{document}
Tags replace the number of a row; see \eqref{eq:star}, \eqref{eq:two} and \ref{eq:raw}.
\begin{align}
a &= b + c \tag{A} \label{eq:star} \\
d &= e + f \\
g &= h \tag*{[iii]} \label{eq:raw} \\
x^2 + y^2 &= z^2 \tag{$\ast\ast$} \\
u &= v \notag \\
p &= q + r \tag{2.1\textit{f}} \\
s &= t \label{eq:two}
\end{align}
and in the starred form, where no row has a number of its own,
\begin{align*}
f(x) &= x^2 \tag{B'} \\
&= x \cdot x \\
&\ge 0 \tag*{Q.E.D.}
\end{align*}
as claimed in \eqref{eq:star}.
\end{document}
